\appendix

% ============================================================================
%  Por que a fase sobrevive à equalização
% ============================================================================
\begin{frame}{Apêndice}{Por que o Equalizador não Remove a Fase}
  \small

  Retomando o espectro do sinal amostrado e aplicando a cadeia de reconstrução:
  \begin{equation*}
    X_p(f) = \frac{\tau}{T_s}\,\mathrm{sinc}(f\tau)\,e^{-j\pi f\tau}
             \sum_{n=-\infty}^{\infty} X(f - nf_s)
  \end{equation*}

  O passa-baixas $H_r(f) = T_s/\tau$ em $|f| \leq f_c$ isola a réplica central
  ($n=0$) e cancela o ganho $\tau/T_s$:
  \begin{equation*}
    X_1(f) = \mathrm{sinc}(f\tau)\,e^{-j\pi f\tau}\,X(f)
  \end{equation*}

  O equalizador $E(f) = 1/\mathrm{sinc}(f\tau)$ é uma função
  \textbf{real e positiva} em $|f| < 1/\tau$, logo cancela o módulo mas deixa a
  fase intacta:
  \begin{equation*}
    X_{rec}(f) = E(f)\,X_1(f) = e^{-j\pi f\tau}\,X(f)
  \end{equation*}

  Pela propriedade do deslocamento no tempo,
  $e^{-j2\pi f t_0}X(f) \longleftrightarrow x(t - t_0)$, com
  $t_0 = \tau/2$:
  \begin{equation*}
    \boxed{\;x_{rec}(t) = x\!\left(t - \frac{\tau}{2}\right)\;}
  \end{equation*}

\end{frame}
